\chapter{Uma transacção de tipo 6, byte a byte}
\label{appendix:pagamento-tipo-seis}
\index{serializacao@serialização!transaccao de tipo 6@transacção de tipo 6}

Tomemos uma transacção de uma entrada e duas saídas, incluída no bloco
\(2\,777\,777\) da lista de blocos de Monero. O seu registo tem 1539
bytes. O nó que o recebe consegue reconstituir dezasseis referências
a saídas anteriores, duas chaves novas, uma taxa, três compromissos de
montante e os dados de duas provas diferentes. Todos partem da mesma
sequência de bytes. Como sabem onde acaba uma coisa e começa a outra?

O exemplar conservado em~\cite{monero-transaccao-tipo-seis} tem o identificador
\begin{center}
\texttt{e89415b95564aa7e3587c91422756ba5}\\
\texttt{303e727996e19c677630309a0d52a7ca}.
\end{center}
As duas linhas são metades de um único resumo de 32 bytes. Voltaremos
a obtê-lo no fim da leitura. Os montantes de Alice e de Bob, no
capítulo~\ref{chapter:transactions}, eram conhecidos do leitor porque
construímos o exemplo a partir das aberturas. Aqui estamos no lugar do nó.
Conhecemos os bytes públicos, incluindo a taxa. As aberturas das saídas
ficaram com os participantes.

\section{A primeira passagem pelo prefixo}

Contaremos posições a partir de zero e escreveremos \([a,b[\) para os
bytes desde \(a\) até \(b-1\). Com esta convenção o comprimento é sempre
\(b-a\), e o limite direito de um campo é o princípio do seguinte. Os
primeiros seis bytes, em hexadecimal, são
\[
                  \mathtt{02\quad00\quad01\quad02\quad00\quad10}.
\]
O primeiro \texttt{02} indica a versão de serialização. O zero na
posição 1 é o tempo de desbloqueio inscrito na transacção. O
\texttt{01} conta uma entrada. Chegámos à posição 3, onde o segundo
\texttt{02} identifica a variante de entrada por chave. O zero
seguinte selecciona a classe de saídas com montante explícito zero.
O \texttt{10} final vale dezasseis em decimal e conta as referências
do anel. É a posição de cada byte na gramática que lhe dá significado.
O mesmo \texttt{02}, lido mais adiante, contará duas saídas.

Para encontrar o fim do anel, o leitor procura o fim de dezasseis
inteiros de comprimento variável. O primeiro ocupa quatro bytes,
\texttt{96 a4 8a 1a}. Nos três primeiros, o bit superior vale um e
anuncia uma continuação. Retirado esse bit, os quatro algarismos em
base 128 são, pela ordem de leitura, \(22,36,10,26\). Temos
\[
 22+36\cdot128+10\cdot16\,384+26\cdot2\,097\,152
       =54\,694\,422.
\]
O segundo, \texttt{b1 b9 bd 01}, dá
\(49+57\cdot128+61\cdot16\,384+2\,097\,152=3\,103\,921\).
O primeiro valor já era um índice global. O segundo é a distância
ao membro seguinte. A segunda consulta do nó será, portanto, ao índice
\(54\,694\,422+3\,103\,921=57\,798\,343\).
A tabela segue esta mesma soma até ao último membro.

\begin{center}\small
\begin{tabular}{r l r r}
Posição no anel & Bytes transmitidos & Distância & Índice consultado\\
\hline
0 & \texttt{96 a4 8a 1a} & 54\,694\,422 & 54\,694\,422\\
1 & \texttt{b1 b9 bd 01} & 3\,103\,921 & 57\,798\,343\\
2 & \texttt{d0 dd ae 02} & 4\,959\,952 & 62\,758\,295\\
3 & \texttt{9f eb 1b} & 456\,095 & 63\,214\,390\\
4 & \texttt{91 de 31} & 814\,865 & 64\,029\,255\\
5 & \texttt{9c b1 1d} & 481\,436 & 64\,510\,691\\
6 & \texttt{dd d2 1b} & 452\,957 & 64\,963\,648\\
7 & \texttt{e6 ea 28} & 669\,030 & 65\,632\,678\\
8 & \texttt{f2 b1 07} & 121\,074 & 65\,753\,752\\
9 & \texttt{d6 61} & 12\,502 & 65\,766\,254\\
10 & \texttt{a0 ab 05} & 87\,456 & 65\,853\,710\\
11 & \texttt{ab d1 04} & 75\,947 & 65\,929\,657\\
12 & \texttt{e2 94 01} & 19\,042 & 65\,948\,699\\
13 & \texttt{96 53} & 10\,646 & 65\,959\,345\\
14 & \texttt{f7 09} & 1\,271 & 65\,960\,616\\
15 & \texttt{f4 86 02} & 33\,652 & 65\,994\,268
\end{tabular}
\end{center}

Os quatro bytes do último índice absoluto seriam necessários para
guardar um número da grandeza de \(65\,994\,268\). Aqui bastaram três
para descrever o último salto. Somando a largura das dezasseis
distâncias obtêm-se 48 bytes. Os dezasseis índices absolutos, todos
maiores que \(128^3-1\) e menores que \(128^4\), ocupariam 64 bytes
na mesma codificação. A economia neste anel foi de 16 bytes.
Há outra, bem maior. Para cada índice, o nó obtém da lista de blocos
uma chave de saída e um compromisso. São \(16\cdot2\cdot32=1024\)
bytes de pontos que as referências permitem ir buscar a dados anteriores.

Por exemplo, ao primeiro índice correspondem a chave de saída que
começa por \texttt{e50f4761} e o compromisso que começa por
\texttt{51e788dd}. Ao último correspondem \texttt{4ee857d0} e
\texttt{a8a9d61d}. Estes pares, conservados com o exemplar, esclarecem
o que há a reconstituir. O nó usa os 32 bytes de cada ponto e
conserva a ordem das linhas. A posição que o signatário pôde abrir
permanece secreta. A proximidade numérica de duas linhas mede a
distância entre os índices de registo.

\section{O que o destinatário pode reconhecer}

À posição 54 começa um ponto com outra função. Os 32 bytes até à
posição 86, exclusive, codificam a imagem de chave \(I\), cujo
prefixo é \texttt{af1f15c6}. A CLSAG volta a usar este mesmo ponto.
A carteira inclui-o uma única vez, no prefixo. Assim pode o nó
referir a autorização que está a verificar à mesma imagem por que
procura no registo de gastos. Se receber o prefixo desacompanhado da
prova e do estado anterior, possui já o nome a procurar, embora lhe
falte a justificação para o registar como gasto.

O byte 86 vale dois. Seguem-se duas saídas de 35 bytes cada. Em
cada uma há um zero para o montante explícito, o discriminante
\texttt{03}, uma chave de uso único de 32 bytes e uma etiqueta de
vista de um byte. As terminações e a passagem de uma saída à outra
podem ler-se sem ambiguidade:
\begin{center}\small
\begin{tabular}{c c l c c}
Saída & Início & Começo da chave & Posição da etiqueta & Etiqueta\\
\hline
0 & 87  & \texttt{700516ab} & 121 & \texttt{52}\\
1 & 122 & \texttt{8d5b9072} & 156 & \texttt{08}
\end{tabular}
\end{center}
A primeira carteira que percorre estas saídas conhece a posição local
que deve empregar na derivação, zero ou um. Pode calcular a etiqueta
esperada para cada posição a partir do segredo partilhado. Uma
coincidência de oito bits leva-a a conferir a chave de uso único
inteira. Isto explica como duas etiquetas diferentes, à vista de todos,
podem prestar um serviço ao destinatário sem publicar o seu endereço.

Na posição 157 lê-se \texttt{2c}, ou 44. É o comprimento do campo
\texttt{extra}. Os seus primeiros 33 bytes são uma etiqueta
\texttt{01} e uma chave pública de transacção de 32 bytes, que começa
por \texttt{ad7b21bd}. Sobram onze. O campo seguinte começa por
\texttt{02 09 01}. A etiqueta exterior \texttt{02} abre um
\emph{nonce} com nove bytes de conteúdo. A etiqueta interior
\texttt{01} dá aos oito restantes o sentido de um identificador de
pagamento cifrado. Esses oito bytes são
\texttt{99 46 ee 95 2c 1e b4 ba}. A conta dos comprimentos fecha:
\[
             (1+32)+(1+1+1+8)=33+11=44.
\]
O próximo byte, na posição 202, pertence já à base RingCT.
Esta leitura diz que dados uma carteira recebeu para tentar uma
derivação. Só uma carteira que a consiga completar e reconhecer a
chave de saída pode atribuir o pagamento a si própria.

\section{Conferir o balanço com os próprios pontos}

O primeiro byte da base RingCT vale seis. Seguem-se os cinco bytes
\texttt{80 b6 b4 f7 16}. Os algarismos da taxa em base 128
são \(0,54,52,119,22\). Recuperamos
\[
  54\cdot128+52\cdot128^2+119\cdot128^3+22\cdot128^4
      =6\,156\,000\,000
\]
unidades atómicas, ou \(0{,}006156\) XMR. Seguem-se dois campos de
oito bytes, um para o montante cifrado de cada saída, e dois pontos
de 32 bytes, um para cada compromisso. Por esta altura, saber a
contagem de saídas do prefixo poupou dois novos contadores. A base
inteira ocupa \(1+5+2\cdot8+2\cdot32=86\) bytes.

Podemos agora efectuar, com dados públicos, a conta que no pagamento
de Alice se fazia com letras. Designemos por \(C_0,C_1\) os
compromissos pela ordem das saídas, por \(f\) a taxa acabada de ler
e por \(\widetilde C\) o pseudo-compromisso da única entrada. Na
tabela seguinte cada codificação de ponto ocupa duas metades
consecutivas de dezasseis bytes. As letras maiúsculas da coluna da
esquerda designam pontos somados pela lei de Edwards.
\begin{center}\small
\begin{tabular}{l l}
Ponto & Codificação, por metades\\
\hline
\(C_0\) & \texttt{0149d07055f4f8fcdde1c12896f8c0ee}\\
 & \texttt{25c12fe389c3d4a4f40c6837163818c2}\\
\(C_1\) & \texttt{49b58fa1fc3331a732e502e043e75a41}\\
 & \texttt{06472192afcb4a6b7544dcc70e914182}\\
\(C_0+C_1\) & \texttt{becf4287ca79a5b50d2fada1fc930b84}\\
 & \texttt{f2e53be1fa9652b691d653fef4b86b3e}\\
\(fH\) & \texttt{5d04bc3408fa1164d72c4514f4eba941}\\
 & \texttt{1ac127d4a849a570d4bc6aecfaf39a05}\\
\(C_0+C_1+fH\) & \texttt{eb83fde4db28cfbb8ec0a785bcf6a725}\\
 & \texttt{3099a7c9424161ae04c12f4bc2146408}
\end{tabular}
\end{center}
A última linha é, byte por byte, a sequência dos últimos 32 bytes
da transacção. Aí se encontra \(\widetilde C\). Para reproduzir a
igualdade, o nó descodifica os dois primeiros pontos, soma-os,
multiplica o gerador de montantes pela taxa e faz a última adição.
Só depois volta a comprimir o resultado. Somar os inteiros escritos
em hexadecimal nas duas linhas seria outra operação.

Este teste teve um resultado inteiramente determinado pelos dados
públicos. Continua a faltar uma informação que mesmo a igualdade
verdadeira deixa em aberto. Quais os inteiros ocultos em \(C_0,C_1\)?
A resposta que a transacção oferece ao nó é uma propriedade de ambos:
cada um cabe em 64 bits. A garantia liga-se exactamente aos dois pontos
da tabela. A assinatura, por sua vez, usará o ponto da última linha
ao tratar os dezasseis pares recuperados no começo. É aqui que as
três verificações, feitas com objectos distintos, voltam a encontrar-se.

\section{Onde começam e acabam as provas}

A base terminou na posição \(202+86=288\). Aí aparece um byte com o
valor um, a contagem das provas de intervalo. Para representar as
duas saídas, a prova usa \(2\cdot64=128\) posições. Sete dobras
levam 128 a uma posição, pelo que a cada um dos dois vectores de
pontos \(L,R\) cabem sete elementos. Esta consequência matemática
pode conferir-se directamente na serialização.
\begin{center}\small
\begin{tabular}{l r r r}
Campo & Início & Fim exclusivo & Comprimento\\
\hline
Contagem das provas & 288 & 289 & 1\\
\(A,A_1,B,r_1,s_1,d_1\) & 289 & 481 & 192\\
Comprimento de \(L\) & 481 & 482 & 1\\
Sete pontos de \(L\) & 482 & 706 & 224\\
Comprimento de \(R\) & 706 & 707 & 1\\
Sete pontos de \(R\) & 707 & 931 & 224\\
Dezasseis respostas da CLSAG & 931 & 1443 & 512\\
Desafio inicial \(c_1\) & 1443 & 1475 & 32\\
Imagem auxiliar codificada \(\bar D\) & 1475 & 1507 & 32\\
Pseudo-compromisso \(\widetilde C\) & 1507 & 1539 & 32
\end{tabular}
\end{center}
Os bytes 481 e 706 valem ambos sete. A entrada única e a lista dos
seus dezasseis candidatos determinam as dimensões seguintes: uma
assinatura e dezasseis respostas. Pelos bytes do desafio inicial pode
o nó começar a percorrer o anel, verificando que, na última posição,
volta ao desafio de partida. A segunda imagem, \(\bar D\), diz respeito
à diferença de compromissos. A imagem principal, usada para reconhecer
o gasto repetido, já estava à posição 54.

O balanço do espaço ocupado é tão exacto como o dos pontos, embora
responda a uma pergunta diferente. Entraram 202 bytes de prefixo e
86 de base. A parte podável recebeu
\[
 \underbrace{1+192+2+448}_{643}
 +\underbrace{512+32+32}_{576}+32=1251
\]
bytes. No total, \(202+86+1251=1539\). A primeira chaveta conta
uma prova com os seus contadores. A segunda conta a CLSAG inteira,
incluindo o desafio. Se se dividisse o comprimento total por 32,
todos os contadores, comprimentos, etiquetas e distâncias do prefixo
se perderiam entre os supostos pontos. A dimensão de um anel
determina quanto custa a sua assinatura em escalares. A localização
concreta dos seus membros ainda intervém no custo em bytes.

\subsection{A taxa por unidade de peso}

Com o comprimento da transacção finalmente em mãos, voltemos aos
cinco bytes que anunciavam a taxa. As duas saídas cabem numa única
prova agregada com capacidade para dois montantes. Neste caso a
parcela de recuperação do peso é zero. O peso é, por conseguinte,
1539. Se o leitor quiser saber a que preço por unidade de peso
ficou este pagamento, pode apurá-lo só com o registo:
\[
           \frac{6\,156\,000\,000}{1539}
                =4\,000\,000\quad\text{unidades atómicas}.
\]
A divisão é exacta. Multiplicando de volta,
\(1500\cdot4\,000\,000=6\,000\,000\,000\) e
\(39\cdot4\,000\,000=156\,000\,000\). Em XMR, as duas parcelas
restituem \(0{,}006+0{,}000156=0{,}006156\). Este quociente
descreve o preço efectivo do exemplar. Para determinar a taxa
mínima admitida numa dada altura seriam ainda necessárias as
regras e os dados dessa altura.

Quanto desse peso coube à prova de intervalo e à autorização da
entrada? Podemos responder mesmo que o mineiro receba uma única
taxa para a transacção inteira. Tomemos as fronteiras da tabela e,
a título de repartição aritmética, multipliquemos cada comprimento
pelo preço efectivo que acabámos de calcular.
\begin{center}\small
\begin{tabular}{l r r}
Parcela & Bytes & Quota da taxa em unidades atómicas\\ \hline
Prefixo & 202 & \(808\,000\,000\)\\
Base RingCT & 86 & \(344\,000\,000\)\\
Prova de intervalo e seus contadores & 643 & \(2\,572\,000\,000\)\\
Assinatura CLSAG & 576 & \(2\,304\,000\,000\)\\
Pseudo-compromisso & 32 & \(128\,000\,000\)\\ \hline
Total & 1539 & \(6\,156\,000\,000\)
\end{tabular}
\end{center}
A prova e a assinatura somam 1219 bytes, aproximadamente \(79{,}2\%\)
do peso. As cinco verbas da direita resultaram todas do mesmo
quociente. O único valor de taxa que a carteira escreveu foi o
total da última linha.

Há uma coincidência que ajuda a perceber a economia da poda neste
exemplar. O nó pode conservar os 202 bytes do prefixo, os 86 da base
e 32 para o resumo da parte podável. Os dados necessários para esta
representação perfazem 320 bytes. A diferença é precisamente
\(1539-320=1219\). O pseudo-compromisso que desaparece tem, por acaso,
o mesmo comprimento que o resumo que fica em seu lugar na conta.
O resumo cobre toda a parte podável. Vejamos exactamente como é
que esses 32 bytes permitem recuperar o identificador publicado.

\section{Dois resumos para duas perguntas}

A pergunta que se faz no momento da assinatura é «que pagamento
autorizam estas respostas?». A pergunta que se faz quando o bloco
enumera transacções é «qual a transacção completa que ocupa este
lugar?». Os dois cálculos partem do mesmo prefixo e da mesma base.
Chamemos \(K(B)\) ao Keccak--256 de uma sequência de bytes \(B\).
Se \(P,B,T\) forem o prefixo, a base e a parte podável, formam-se
\[
       h_P=K(P),\qquad h_B=K(B),\qquad h_T=K(T).
\]
A concatenação conserva os 32 bytes de cada valor pela ordem
indicada. Todas as contas abaixo usam os bytes das codificações,
sem os converter nos caracteres com que os imprimimos.

Para a mensagem da CLSAG, substitui-se o terceiro resumo por
\(h_\pi\), calculado a partir dos elementos da prova de intervalo.
A sequência que entra em \(h_\pi\) tem \(192+224+224=640\) bytes.
São os três intervalos \([289,481[\), \([482,706[\) e \([707,931[\),
colados por esta ordem. O contador das provas e os comprimentos dos
dois vectores já cumpriram a sua função durante a leitura. A ordem
dos elementos permanece exactamente a da prova.
\begin{center}\small
\begin{tabular}{l l}
Resumo & Duas metades do resultado\\
\hline
\(h_P\) & \texttt{02465363e123a1d1eeab20d6cec5619d}\\
 & \texttt{18b483a1cc7baa9489f4367a946b3bca}\\
\(h_B\) & \texttt{2bccf4bd1fa784d7965944eabbd8b286}\\
 & \texttt{b6ebfbbaebb7eb10c82e7c486da66026}\\
\(h_T\) & \texttt{824cebe4ce6fb25648d92821c101eea7}\\
 & \texttt{4d81dc825a0d187d883698b712028b3f}\\
\(h_\pi\) & \texttt{7f7243a4c24f712626f8ed8189758c7d}\\
 & \texttt{04408cdfc64c670084f026bd52850792}
\end{tabular}
\end{center}
O nó calcula a mensagem \(\mathfrak m=K(h_P\|h_B\|h_\pi)\) para
verificar as respostas da única entrada. Obtém
\begin{center}
\texttt{c5a76ffd7d8916210332afa07a253368}\\
\texttt{f7e1eea1f385f393a1f3fa8b6d1f2f4b}.
\end{center}
No cálculo do identificador, \(K(h_P\|h_B\|h_T)\), o último resumo
abrange as respostas da própria assinatura, a prova com os seus
comprimentos e o pseudo-compromisso final. O resultado começa por
\texttt{e89415b9} e termina em \texttt{0d52a7ca}: é o identificador
reproduzido na abertura do apêndice. A mensagem assinada pôde ser
calculada antes de existirem as respostas. O identificador permite
reconhecer também essas respostas depois de reunidas.

Imaginemos agora o mesmo registo num nó que já podou os dados desta
transacção. Ficaram disponíveis o prefixo, a base e \(h_T\). O nó
ainda pode recompor o identificador inteiro. Um leitor que queira
conferir de novo a igualdade dos três pontos precisa dos bytes de
\(\widetilde C\), que pertenciam à parte podável. Para verificar a
CLSAG precisa ainda da assinatura e dos pares de pontos do anel.
Quem lhe entregar os 1251 bytes pode ser confrontado com \(h_T\).
A correspondência do resumo identifica os dados procurados. A
verificação algébrica dará então a esses dados o seu significado
como autorização de um gasto concreto.
